\documentclass[10pt,a4paper]{amsart}
\usepackage{amsmath,amssymb,amsthm,graphicx}
% Retained preamble for the Stacks Project chapter divisors.tex (Divisors).
% The chapter file pinned-original-source.tex carries only \input{preamble} on its first line, so no
% preamble text lives inside the chapter itself.  The Stacks Project uses one single shared preamble for
% every chapter; the snapshot stages it verbatim in sources/stacks-chow/preamble.tex.  The block below is
% copied mechanically, line for line, from that file's lines 14-18 (the Xy-pic setup) and lines 154-217
% (the native theorem declarations together with the mathematical macros).  Nothing was added, removed or
% redefined.  The wrapper class and the amsmath/amssymb/amsthm/graphicx base packages are supplied by the
% assembler, so the class/\externaldocument part of the shared preamble is not repeated here.
% For commutative diagrams we use Xy-pic
\usepackage[all]{xy}

% We use 2cell for 2-commutative diagrams.
\xyoption{2cell}
% Theorem environments.
%
\theoremstyle{plain}
\newtheorem{theorem}[subsection]{Theorem}
\newtheorem{proposition}[subsection]{Proposition}
\newtheorem{lemma}[subsection]{Lemma}

\theoremstyle{definition}
\newtheorem{definition}[subsection]{Definition}
\newtheorem{example}[subsection]{Example}
\newtheorem{exercise}[subsection]{Exercise}
\newtheorem{situation}[subsection]{Situation}

\theoremstyle{remark}
\newtheorem{remark}[subsection]{Remark}
\newtheorem{remarks}[subsection]{Remarks}

\numberwithin{equation}{subsection}

% Macros
%
\def\lim{\mathop{\mathrm{lim}}\nolimits}
\def\colim{\mathop{\mathrm{colim}}\nolimits}
\def\Spec{\mathop{\mathrm{Spec}}}
\def\Hom{\mathop{\mathrm{Hom}}\nolimits}
\def\Ext{\mathop{\mathrm{Ext}}\nolimits}
\def\SheafHom{\mathop{\mathcal{H}\!\mathit{om}}\nolimits}
\def\SheafExt{\mathop{\mathcal{E}\!\mathit{xt}}\nolimits}
\def\Sch{\mathit{Sch}}
\def\Mor{\mathop{\mathrm{Mor}}\nolimits}
\def\Ob{\mathop{\mathrm{Ob}}\nolimits}
\def\Sh{\mathop{\mathit{Sh}}\nolimits}
\def\NL{\mathop{N\!L}\nolimits}
\def\CH{\mathop{\mathrm{CH}}\nolimits}
\def\proetale{{pro\text{-}\acute{e}tale}}
\def\etale{{\acute{e}tale}}
\def\QCoh{\mathit{QCoh}}
\def\Ker{\mathop{\mathrm{Ker}}}
\def\Im{\mathop{\mathrm{Im}}}
\def\Coker{\mathop{\mathrm{Coker}}}
\def\Coim{\mathop{\mathrm{Coim}}}

% Boxtimes
%
\DeclareMathSymbol{\boxtimes}{\mathbin}{AMSa}{"02}

%
% Macros for moduli stacks/spaces
%
\def\QCohstack{\mathcal{QC}\!\mathit{oh}}
\def\Cohstack{\mathcal{C}\!\mathit{oh}}
\def\Spacesstack{\mathcal{S}\!\mathit{paces}}
\def\Quotfunctor{\mathrm{Quot}}
\def\Hilbfunctor{\mathrm{Hilb}}
\def\Curvesstack{\mathcal{C}\!\mathit{urves}}
\def\Polarizedstack{\mathcal{P}\!\mathit{olarized}}
\def\Complexesstack{\mathcal{C}\!\mathit{omplexes}}
% \Pic is the operator that assigns to X its picard group, usage \Pic(X)
% \Picardstack_{X/B} denotes the Picard stack of X over B
% \Picardfunctor_{X/B} denotes the Picard functor of X over B
\def\Pic{\mathop{\mathrm{Pic}}\nolimits}
\def\Picardstack{\mathcal{P}\!\mathit{ic}}
\def\Picardfunctor{\mathrm{Pic}}
\def\Deformationcategory{\mathcal{D}\!\mathit{ef}}

% --- appended 2026-10-12 (Chief): the project own definitions of the reference/slogan/history
% display environments, copied mechanically from the shared Stacks preamble lines 8-12
% (\usepackage{verbatim} plus the three \newenvironment lines, which the chapter preamble keeps
% verbatim from the project preamble). Nothing added, removed or redefined: these environments
% carry no mathematical content and the project itself typesets them as comments. Without these
% three lines a counted statement whose own text uses them fails with Environment slogan undefined.
\usepackage{verbatim}
\newenvironment{reference}{\comment}{\endcomment}
\newenvironment{slogan}{\comment}{\endcomment}
\newenvironment{history}{\comment}{\endcomment}
\usepackage{cleveref}
\crefname{originallemma}{original lemma}{original lemmas}
\Crefname{originallemma}{Original lemma}{Original lemmas}
\makeatletter
\newlabel{algebra-lemma-ass-zero-divisors}{{00LD}{source}{}{sourcepointer-1}{}}
\newlabel{algebra-lemma-ass-zero-divisors@cref}{{[originallemma][0][]00LD}{[1][0][]source}{}{}{}}
\makeatother
\crefname{originallemma}{original lemma}{original lemmas}
\Crefname{originallemma}{Original lemma}{Original lemmas}
\makeatletter
\newlabel{algebra-lemma-associated-primes-localize}{{0310}{source}{}{sourcepointer-2}{}}
\newlabel{algebra-lemma-associated-primes-localize@cref}{{[originallemma][0][]0310}{[1][0][]source}{}{}{}}
\makeatother
\crefname{originallemma}{original lemma}{original lemmas}
\Crefname{originallemma}{Original lemma}{Original lemmas}
\makeatletter
\newlabel{algebra-lemma-silly}{{00DS}{source}{}{sourcepointer-3}{}}
\newlabel{algebra-lemma-silly@cref}{{[originallemma][0][]00DS}{[1][0][]source}{}{}{}}
\makeatother
\crefname{originallemma}{original lemma}{original lemmas}
\Crefname{originallemma}{Original lemma}{Original lemmas}
\makeatletter
\newlabel{lemma-meromorphic-section-restricts-to-zero}{{08I7}{source}{}{sourcepointer-4}{}}
\newlabel{lemma-meromorphic-section-restricts-to-zero@cref}{{[originallemma][0][]08I7}{[1][0][]source}{}{}{}}
\makeatother
\title{The sheaf of meromorphic functions on a locally Noetherian scheme}
\author{The Stacks Project Authors}
\date{stacks-project:divisors.tex:local-snapshot-20261012; GNU FDL 1.2 or later}
\begin{document}
\maketitle
\setcounter{subsection}{1}
\noindent\textit{Editorial note.} All delivered mathematical text is the author's own, copied from the pinned Stacks Project chapter file divisors.tex; no statement or proof step was rewritten, reordered or removed.\par\medskip
\noindent\textit{Editorial note.} Carrier window 3: statement and proof both lie inside the window 6000-8000 granted to this lane by the owner lane lane-s13 on 2026-10-14.\par\medskip
\noindent\textit{Editorial note.} The references of the counted spans are recorded as non-load-bearing navigation pointers; the argument that consumes them stays inside the counted proof span.\par\medskip
\noindent\textit{Editorial note.} Preamble: sources/stacks-divisors/retained-preamble.tex (88 lines, verbatim present).\par\medskip
\noindent\textit{Allowed standard prerequisites (editorial).}\begin{itemize}
\item A meromorphic section vanishing at a point vanishes on a neighbourhood. Lemma lemma-meromorphic-section-restricts-to-zero, used at proof lines 6300-6302. Produces the element e with e\textasciicircum{}n(ad - cf) = 0.
\item Localisation and associated primes. Algebra, Lemma algebra-lemma-associated-primes-localize, used at proof lines 6306-6310. Shows I is not contained in any associated prime.
\item An ideal not contained in the union of the associated primes has a nonzerodivisor. Algebra, Lemmas algebra-lemma-silly and algebra-lemma-ass-zero-divisors, used at proof lines 6310-6313. Produces the nonzerodivisor z in I.
\end{itemize}\medskip
\ifdefined\hypertarget\hypertarget{sourcepointer-1}{}\fi
\noindent\textit{Source reference.} Original lemma 00LD: focused-seeds/AGP/sources/stacks-snapshot-20261011/tags.txt, source lines 783--783. Navigation pointer into the Stacks Project; the standard fact it supplies is applied inside the counted proof span and the argument resting on it stays there.\par\medskip
\ifdefined\hypertarget\hypertarget{sourcepointer-2}{}\fi
\noindent\textit{Source reference.} Original lemma 0310: focused-seeds/AGP/sources/stacks-snapshot-20261011/tags.txt, source lines 3967--3967. Navigation pointer into the Stacks Project; the standard fact it supplies is applied inside the counted proof span and the argument resting on it stays there.\par\medskip
\ifdefined\hypertarget\hypertarget{sourcepointer-3}{}\fi
\noindent\textit{Source reference.} Original lemma 00DS: focused-seeds/AGP/sources/stacks-snapshot-20261011/tags.txt, source lines 508--508. Navigation pointer into the Stacks Project; the standard fact it supplies is applied inside the counted proof span and the argument resting on it stays there.\par\medskip
\ifdefined\hypertarget\hypertarget{sourcepointer-4}{}\fi
\noindent\textit{Source reference.} Original lemma 08I7: focused-seeds/AGP/sources/stacks-snapshot-20261011/tags.txt, source lines 10805--10805. Navigation pointer into the Stacks Project; the standard fact it supplies is applied inside the counted proof span and the argument resting on it stays there.\par\medskip
\begin{lemma}
\label{lemma-locally-Noetherian-K}
Let $X$ be a locally Noetherian scheme.
\begin{enumerate}
\item For any $x \in X$ we have $\mathcal{S}_x \subset \mathcal{O}_{X, x}$
is the set of nonzerodivisors, and hence $\mathcal{K}_{X, x}$
is the total quotient ring of $\mathcal{O}_{X, x}$.
\item For any affine open $U \subset X$ the ring
$\mathcal{K}_X(U)$ equals the total quotient ring of $\mathcal{O}_X(U)$.
\end{enumerate}
\end{lemma}

\begin{proof}
To prove this lemma we may assume $X$ is the spectrum of a Noetherian
ring $A$. Say $x \in X$ corresponds to $\mathfrak p \subset A$.

\medskip\noindent
Proof of (1). It is clear that $\mathcal{S}_x$ is contained
in the set of nonzerodivisors of $\mathcal{O}_{X, x} = A_\mathfrak p$.
For the converse, let $f, g \in A$, $g \not \in \mathfrak p$ and
assume $f/g$ is a nonzerodivisor in $A_{\mathfrak p}$. Let
$I = \{a \in A \mid af = 0\}$. Then we see that $I_{\mathfrak p} = 0$ by
exactness of localization. Since $A$ is Noetherian we see that $I$
is finitely generated and hence that $g'I = 0$ for some $g' \in A$,
$g' \not \in \mathfrak p$. Hence $f$ is a nonzerodivisor
in $A_{g'}$, i.e., in a Zariski open neighbourhood of $\mathfrak p$.
Thus $f/g$ is an element of $\mathcal{S}_x$.

\medskip\noindent
Proof of (2). Let $f \in \Gamma(X, \mathcal{K}_X)$ be a meromorphic function.
Set $I = \{a \in A \mid af \in A\}$. Fix a prime $\mathfrak p \subset A$
corresponding to the point $x \in X$. By (1) we can write the image of $f$
in the stalk at $\mathfrak p$ as $a/b$, $a, b \in A_{\mathfrak p}$ with
$b \in A_{\mathfrak p}$ not a zerodivisor. Write $b = c/d$ with
$c, d \in A$, $d \not \in \mathfrak p$. Then $ad - cf$ is a section of
$\mathcal{K}_X$ which vanishes in an open neighbourhood of $x$. Say it
vanishes on $D(e)$ with $e \in A$, $e \not \in \mathfrak p$. Then
$e^n(ad - cf) = 0$ for some $n \gg 0$ by
Lemma \ref{lemma-meromorphic-section-restricts-to-zero}.
Thus $e^nc \in I$ and $e^nc$ maps to a nonzerodivisor in
$A_{\mathfrak p}$. Let
$\text{Ass}(A) = \{\mathfrak q_1, \ldots, \mathfrak q_t\}$ be the
associated primes of $A$. By looking at $IA_{\mathfrak q_i}$ and
using Algebra, Lemma \ref{algebra-lemma-associated-primes-localize}
the above says that
$I \not \subset \mathfrak q_i$ for each $i$. By
Algebra, Lemma \ref{algebra-lemma-silly}
there exists an element $z \in I$, $z \not \in \bigcup \mathfrak q_i$.
By Algebra, Lemma \ref{algebra-lemma-ass-zero-divisors}
we see that $z$ is not a zerodivisor on $A$.
Hence $f = (zf)/z$ is an element of the total ring of fractions of $A$.
This proves (2).
\end{proof}

\end{document}
